\section{The full order model}\label{app:formal-model}

A rule is $r_i=(a_i,c_i)$ with logical meaning $\phi_i=a_i\Rightarrow c_i$; $\Phi$ is the conjunction of the retained rule formulas. A valuation assigns every symbol one value in a single snapshot.
For the order example, we use the following abbreviations alongside the identifiers saved in the formal model.

\begin{table}[H]
\centering
\begin{tabularx}{\linewidth}{@{}clX@{}}
\toprule
Symbol & Saved identifier & Meaning \\
\midrule
$s$ & \code{order_submitted} & A submission event occurs. \\
$v$ & \code{inventory_available} & Stock is available. \\
$f$ & \code{order_fulfilled} & The order is fulfilled. \\
$b$ & \code{order_on_backorder} & The order is backordered. \\
$c$ & \code{order_canceled} & A cancellation event occurs. \\
$q$ & \code{order_canceled_and_refunded} & The canceled-and-refunded state holds. \\
$u$ & \code{payments_refunded} & The payments are refunded. \\
\bottomrule
\end{tabularx}
\caption{Stable abbreviations for the order worked example. All seven symbols are Boolean.}
\end{table}

An atomic rule is $r_i=(a_i,c_i)$, with antecedent $a_i$ and consequent $c_i$. Its logical meaning is $\phi_i=a_i\Rightarrow c_i$. The subscripted $c_i$ means a rule's consequent; the standalone $c$ in the order example is its cancellation event. With $\Phi=\bigwedge_i\phi_i$, the five rules become
\begin{align*}
\phi_1&=(s\land v)\Rightarrow f,&
\phi_2&=(s\land\neg v)\Rightarrow b,\\
\phi_3&=\True\Rightarrow\neg b,&
\phi_4&=c\Rightarrow u,\\
\phi_5&=q\Rightarrow\neg f.&&
\end{align*}

The persisted formulas use SMT-LIB syntax~\citep{smtlib}. For example, R2 is written as an antecedent \code{(and order_submitted (not inventory_available))} and a consequent \code{order_on_backorder}. The vocabulary object parses these strings and constructs Z3 expressions. The checker is not interpreting arbitrary Python generated by the language model.
Appendix~\ref{app:walkthrough} shows how the saved requirement, symbol, translation and finding correspond.

\source{speccheck/logic.py; speccheck/checks.py; specs/order/formal.json}

\subsection{Solver queries and assumptions}

Write $\SAT(F)$ when the solver returns a satisfying assignment for $F$, and $\UNSAT(F)$ when it establishes that no such assignment exists. The implementation also allows an \emph{unknown} answer, including on timeout. Unknown establishes neither result. Each solver query has a five-second timeout; the sensitivity annotations discussed later need an additional qualification.

We repeatedly use three patterns:
\begin{align}
\text{Equivalence: }&\quad \UNSAT(\phi\oplus\psi),\label{eq:equivalence}\\
\text{Entailment: }&\quad \UNSAT(B\land\neg d),\label{eq:entailment}\\
\text{Incompatibility in a context: }&\quad \UNSAT(T\land\Phi).\label{eq:incompatibility}
\end{align}
Here $B$ is a conjunction of assumptions, $d$ is a candidate consequence, and $T$ describes a triggering context. If the equivalence query is satisfiable, its assignment shows how the two readings differ. If the incompatibility query is unsatisfiable, we still need to check that $T$ itself is feasible before calling the result a conflict.

\example{Why simply checking all five rules is insufficient.}
Set $s=c=q=\False$ and $b=\False$. The conditional rules impose no demands, R3 is obeyed, and $\Phi$ is satisfiable. That does not resolve the backorder problem. Set $T=s\land\neg v$ instead. Now R2 forces $b$ and R3 forces $\neg b$. We have $\SAT(T)$ but $\UNSAT(T\land\Phi)$. The triggering context is what turns a vague global consistency question into the concrete problem we wanted to expose.

\subsection{Assumptions are not requirements}

Let $B$ hold relations proposed by the language model, such as a possible connection between removal and retention. These formulas can help reveal conflicts that the chosen symbols otherwise hide. They can also be wrong. The implementation first searches for conflicts using the rules alone, then searches with assumptions and labels any additional findings that depend on them. Assumptions may cast doubt on a rules-only finding; they must not silently remove it.

\example{An assumption can hide a clash between requirements.}
Suppose an assumed fact says that fulfillment and the canceled-and-refunded state never coexist: $\neg(f\land q)$. The trigger $s\land v\land q$ is still feasible. But under that assumption, R1 alone becomes impossible to obey there: it demands $f$ while the assumption forbids it. An assumptions-only pass would classify the case as a situation forbidden by one rule, rather than report the R1/R5 clash. Running the requirements-only check first preserves that clash for the author to inspect. The policy question stays visible.

Z3 checks the formulas it receives. To assess a finding, we need to know which formulas came from translated requirements, which came from assumed relations, and which interpretation choices affect the conclusion.
\source{speccheck/logic.py; speccheck/checks.py; specs/order/formal.json}

